\chapter{The Naming Dictionary}\label{sec:TheNamingDictionary}
The names of program elements provide an implicit contract (or at least a kind of commitment) between the original author of the program and its users/maintainers. But because people understand words differently, I have added a dictionary of common verbs and their implicit contracts here, together with a list of suffixes for class names.

%------------------------------------------------------------------------------

\section{Verbs}\label{sec:TheNamingDictionary:Verbs}
This chapter provides a list of verbs\footnote{Ok, some names or prefixes are not verbs, like \bsq{main}, \bsq{new} \bsq{from} or \bsq{to}~…} to be used with method names and describes their implicit contracts. These verbs are usually prefixes to a method name, although some of them could be used as standalone names, too. The form that used more often is mentioned first.

The list is by far not exhaustive, and it is in no way authoritative. You may find good reasons to use the names in a manner different from the suggested; one may be that your project has already established it differently.

Feel free to amend and enhance the table according to your needs! This is just a starting point. Look for the file
\begin{itemize}
	\item{\verb#naming/DomainGlossary.tbl.tex#}
\end{itemize}
in the \TeX/\LaTeX\index{TeX}\index{LaTeX} sources for this book on GitHub\index{GitHub} \autocite{TQUADRAT_ORG_DOCUMENT_REPOSITORY} to place your modifications.

If you make the dictionary a part of your coding conventions, you should add a comment to methods\index{method} that behave differently from what is expected based on its name. On the other hand, that a method\index{method} behaves according to the \bsq{contract} does not free you from providing a proper Javadoc comment that describes the purpose of the method\index{method} in detail, together with the parameters, return values and exceptions, as described in \tqfullref{sec:DocumentationComments} and more specific in \tqfullvref{sec:MethodComment}.

%------------------------------------------------------------------------------

\renewcommand{\cellalign}{tl}
\LTXtable{\linewidth}{naming/Verbs.tbl.tex}

%------------------------------------------------------------------------------

\section{Suffixes for Class Names}\label{sec:SuffixesForClassNames}
This chapter lists defined suffixes for class names and their function.

\renewcommand{\cellalign}{tl}
\LTXtable{\linewidth}{naming/ClassNameSuffixes.tbl.tex}

\section{Reserved Names}

The table below lists names for program elements that are somehow reserved. That means that you should avoid to use these names in another context or meaning as described in this document.

The list is definitely not exhaustive, and it is also not authoritative, but a strong recommendation.

\renewcommand{\cellalign}{tl}
\LTXtable{\linewidth}{naming/Reserved.tbl.tex}

%------------------------------------------------------------------------------

\section{Domain Vocabulary}\label{sec:DomainVocabulary}
If your project was designed using \bdq{Domain-Driven Design}\index{domain-driven design} \autocite{Evans:DomainDrivenDesign}, and you will adjust this Coding Conventions for your project, you should consider to add your project's vocabulary here.

Such a glossary may look like the sample structure below that consist from a table that lists the terms together with a brief description, serving as a kind of index to the terms, and a section that describes the terms in details. If that section gets larger, you may consider to move it to the \nameref{sec:Appendices}.

If you want to add your stuff here, look for the files
\begin{itemize}
\item{\verb#naming/DomainGlossary.tex#}
\item{\verb#naming/DomainGlossary.tbl.tex#}
\end{itemize} 
in the \TeX/\LaTeX\index{TeX}\index{LaTeX} sources for this book on GitHub\index{GitHub} \autocite{TQUADRAT_ORG_DOCUMENT_REPOSITORY}.

%------------------------------------------------------------------------------

\clearpage
\subsection{The Ubiquitous Language Glossary} 
\renewcommand{\cellalign}{tl}
\LTXtable{\linewidth}{naming/DomainGlossary.tbl.tex}

\subsubsection{The Terms}

% The suggested structure for a Domain Glossary

% If you avoid references to other parts of the Coding Conventions document, you can use this
% file also as part of a standalone Domain Glossary, using DomainGlossary.tbl.tex as the
% index to it.

\noindent\rule{\textwidth}{0.4pt}

\paragraph{[\textit{Term Name}]}\label{ddterm:TermName}

\subparagraph{Bounded Context:} [Which context this definition applies to]

\subparagraph{Definition:} 
[Clear, precise definition as agreed upon by domain experts and dev team]

\subparagraph{Synonyms/Aliases:} 
[Other names used for this concept, if any – flag if inconsistent]

\subparagraph{Type:} 
[Entity | Value Object | Aggregate | Aggregate Root | Domain Event | Domain Service | Policy | Specification | Enum/Category]

\subparagraph{Attributes/Properties:} 
[Key attributes, if it's a model element]

\subparagraph{Related Terms:} 
[Links to other glossary entries this term relates to or depends on]

\subparagraph{Example:} 
[Concrete example in plain business language]

\subparagraph{Code Reference:} 
[Class/interface name in the codebase – keeps code and language in sync]

\subparagraph{Notes/Constraints:} 
[Business rules, invariants, edge cases]

\noindent\rule{\textwidth}{0.4pt}

%==============================================================================
% End of file!!
%==============================================================================


%==============================================================================
% End of file!!
%==============================================================================
